\chapter{Cycling}

\section*{Definition and Overview}
Cycling is riding a bicycle for transport, recreation, exercise, or sport. It is a low-impact, accessible form of physical activity that can be enjoyed by people of almost any age and fitness level. Regular cycling improves cardiovascular fitness, strengthens muscles, helps control weight, and supports mental wellbeing, and cycling for transport also reduces traffic congestion and emissions. It does carry risks, principally injuries from falls and collisions, which are reduced by helmets, visibility, road skills, and safe infrastructure. This chapter covers the health benefits of cycling, how to cycle safely, and common cycling-related injuries.

\section*{Epidemiology}
Cycling is among the most popular forms of exercise and transport worldwide. In Australia, a substantial proportion of adults and children ride regularly, and cycling for recreation and commuting has grown. Deaths and serious injuries from cycling have fallen over recent decades, and helmets, which are compulsory in Australia, have made a major contribution. Cycling injuries still occur, with falls being the most common mechanism and head injuries the most serious risk. The health benefits of cycling greatly outweigh the risks for the vast majority of people.

\section*{Aetiology and Pathophysiology}
The health effects of cycling come from regular physical activity, while injuries result from falls and collisions.

\begin{itemize}
    \item \textbf{Cardiovascular fitness:} cycling is an aerobic exercise that strengthens the heart and lungs
    \item \textbf{Muscle strength:} it works the legs, glutes, and core
    \item \textbf{Weight control:} regular riding burns energy and supports a healthy weight
    \item \textbf{Mental health:} exercise releases endorphins and reduces stress and anxiety
    \item \textbf{Bone and joint health:} it is low-impact and suits people with joint problems
    \item \textbf{Injury mechanisms:} falls from the bike, collisions with vehicles, and overuse injuries
    \item \textbf{Head injury:} the most serious risk, reduced by wearing an approved helmet
\end{itemize}

\section*{Clinical Features}
Cycling itself has no symptoms; the relevant features are the benefits and the injuries.

\begin{itemize}
    \item Improved fitness, stamina, and energy levels with regular riding
    \item Weight loss or maintenance when combined with a healthy diet
    \item Reduced stress and improved mood
    \item Injuries from falls: abrasions, bruises, sprains, and fractures
    \item Head injury, which can be serious without a helmet
    \item Overuse injuries: knee pain, saddle soreness, and hand or wrist discomfort
    \item Discomfort from an incorrectly fitted bicycle
\end{itemize}

\section*{Differential Diagnosis}
Before starting or increasing cycling, especially with new symptoms, consider the cause.

\begin{itemize}
    \item Chest pain on exertion, which needs medical assessment rather than being assumed to be fitness-related
    \item Joint pain from injury versus overuse or arthritis
    \item Shortness of breath from deconditioning versus heart or lung disease
    \item Dizziness during exercise, which may indicate a medical problem
    \item Unexplained fatigue, which may have a medical cause
\end{itemize}

\section*{When to Seek Medical Care}
See a doctor if you have not been active and plan to take up vigorous cycling, if you have heart disease, high blood pressure, or other chronic conditions, or if exercise causes chest pain, severe breathlessness, or dizziness.

\textbf{Call 000 (or local emergency services) for:}
\begin{itemize}
    \item A cycling accident with a suspected head, neck, or spine injury
    \item Loss of consciousness, confusion, or vomiting after a fall
    \item Severe bleeding or a suspected fracture
    \item Chest pain, difficulty breathing, or collapse during or after riding
    \item A collision with a vehicle, even if you feel well at first
\end{itemize}

\section*{Diagnosis and Investigations}
Assessment is mainly a matter of good preparation, with medical review when needed.

\begin{itemize}
    \item \textbf{Pre-exercise review:} a health check for people with chronic conditions or new symptoms
    \item \textbf{Bike fit:} adjustment of seat height, handlebars, and posture to prevent discomfort
    \item \textbf{Injury assessment:} examination and imaging such as X-ray for suspected fractures
    \item \textbf{Head injury assessment:} urgent evaluation after any fall involving the head
    \item \textbf{Medical tests:} ECG or other tests if exercise-related symptoms occur
\end{itemize}

\section*{Management}
Getting the most from cycling is about training, safety, and equipment.

\begin{itemize}
    \item \textbf{Progressive training:} start gently and increase distance and intensity gradually
    \item \textbf{Helmet:} wear an approved, correctly fitted helmet on every ride
    \item \textbf{Visibility:} lights, bright clothing, and reflective gear
    \item \textbf{Road safety:} follow road rules, use bike lanes, and signal turns
    \item \textbf{Bike maintenance:} check brakes, tyres, and chain regularly
    \item \textbf{Injury treatment:} rest, ice, and first aid for minor injuries; medical care for serious ones
    \item \textbf{Recovery and rehabilitation:} physiotherapy for overuse injuries
    \item \textbf{Combined activity:} mix cycling with strength and flexibility exercise
\end{itemize}

\section*{Complications}
\begin{itemize}
    \item Falls causing abrasions, fractures, and head injuries
    \item Collisions with vehicles or pedestrians
    \item Head injury with concussion or more serious brain injury
    \item Overuse injuries such as knee pain and saddle sores
    \item Sun exposure and dehydration on long rides
    \item Safety risks at night or in poor weather
    \item The small risk of a cardiac event during vigorous exercise in people with undiagnosed heart disease
\end{itemize}

\section*{Prognosis and Outcomes}
Cycling is one of the healthiest forms of exercise, and regular riders have lower rates of heart disease, diabetes, obesity, and early death. Most cycling injuries are minor and heal quickly, and serious injuries are substantially reduced by helmets and safe riding. With sensible progression, appropriate equipment, and road safety, cycling can be continued safely into older age and delivers lifelong health benefits.

\section*{Prevention and Self-Care}
\begin{itemize}
    \item Wear a correctly fitted approved helmet on every ride
    \item Be visible: lights, reflective clothing, and predictable riding
    \item Follow the road rules and ride defensively
    \item Check the bike before riding: brakes, tyres, and chain
    \item Build up gradually and listen to your body
    \item Stay hydrated and protect against the sun
    \item Have a medical check before starting vigorous exercise if you have risk factors
    \item Treat minor injuries promptly and seek review for persistent pain
\end{itemize}

\section*{Sources and Further Reading}
The following reputable sources provide further information for healthcare students and patients seeking reliable, current advice about cycling and cycling safety.

\begin{itemize}
    \item Bicycle Network. \textit{Cycling advice and safety.} https://www.bicyclenetwork.com.au/
    \item Transport for NSW. \textit{Cycling safety and road rules.} https://www.transport.nsw.gov.au/
    \item Healthdirect Australia. \textit{Physical activity and cycling.} https://www.healthdirect.gov.au/
    \item Better Health Channel (Victoria). \textit{Cycling for fitness.} https://www.betterhealth.vic.gov.au/
    \item Cycling Australia. \textit{Sport cycling information.} https://www.cycling.org.au/
    \item World Health Organization. \textit{Physical activity fact sheet.} https://www.who.int/
\end{itemize}
